\begin{explanationbox}
	\textbf{\textcolor{CExplanation}{一句话直觉}}

	椭圆上的点到直线的距离问题，可以通过椭圆的参数方程转化为三角函数的值域问题，从而得到统一的表达式。

	\medskip

	\textbf{\textcolor{CExplanation}{核心拆解}}

	\begin{description}[style=nextline, leftmargin=2.8em, labelsep=0.5em, itemsep=0.45em]
		\item[\textbf{要点一（参数化转化）：}]
		      设椭圆上点为
		      \[
			      (a\cos\theta,\,b\sin\theta),
		      \]
		      代入点到直线距离公式，得
		      \[
			      d=
			      \frac{|Aa\cos\theta+Bb\sin\theta+C|}{\sqrt{A^2+B^2}}.
		      \]

		\item[\textbf{要点二（三角辅助角）：}]
		      设
		      \[
			      M=Aa\cos\theta+Bb\sin\theta,
		      \]
		      利用辅助角公式知 \(M\) 的取值范围是
		      \[
			      [-R,R],
		      \]
		      其中
		      \[
			      R=\sqrt{a^2A^2+b^2B^2}.
		      \]

		\item[\textbf{要点三（分情况讨论）：}]
		      对 \(|M+C|\) 的最值：

		      若 \(|C|>R\)，则
		      \[
			      |M+C|_{\min}=|C|-R,\qquad |M+C|_{\max}=|C|+R.
		      \]

		      若 \(|C|\le R\)，则
		      \[
			      |M+C|_{\min}=0,\qquad |M+C|_{\max}=|C|+R.
		      \]

		      分别除以 \(\sqrt{A^2+B^2}\)，即得距离的最小值与最大值。
	\end{description}

	\medskip

	\textbf{\textcolor{CExplanation}{第二视角（直线向椭圆靠近：平行支撑切线）}}

	前一种方法是\textbf{让点在椭圆上运动}，不断计算该点到固定直线的距离。
	另一种想法是\textbf{固定原直线的方向，沿法向量平移一条平行直线}：
	最外侧刚好碰到椭圆的位置，必定是一条切线。这两条位于椭圆两侧的平行切线，
	恰好夹住整条椭圆。
	这里的“从直线看椭圆”是换一种观察方式；	extbf{原问题仍是椭圆上的动点到固定直线的距离}，
	不是让动点在无限延长的直线上随意运动。

	设椭圆上切点为 \(P(x_0,y_0)\)，该点的切线是
	\[
		\frac{x_0x}{a^2}+\frac{y_0y}{b^2}=1.
	\]
	它的法向量为
	\[
		\vec n_P=\left(\frac{x_0}{a^2},\frac{y_0}{b^2}\right).
	\]
	若切线与已知直线 \(Ax+By+C=0\) 平行，则二者法向量成比例：
	\[
		\left(\frac{x_0}{a^2},\frac{y_0}{b^2}\right)
		=\lambda(A,B).
	\]
	与椭圆方程联立，可得到两个相对的切点
	\[
		P_\pm=\pm\left(\frac{a^2A}{R},\frac{b^2B}{R}\right),
	\]
	以及两条\textbf{平行支撑切线}
	\[
		l_+:Ax+By-R=0,\qquad l_-:Ax+By+R=0.
	\]
	这意味着椭圆全部位于
	\[
		-R\le Ax+By\le R
	\]
	这条由两切线围成的带状区域内。

	已知直线 \(l\) 到这两条切线的距离分别为
	\[
		D_+=\frac{|C+R|}{\sqrt{A^2+B^2}},\qquad
		D_-=\frac{|C-R|}{\sqrt{A^2+B^2}}.
	\]
	\textbf{最远距离}是 \(D_+,D_-\) 中较大的一个；
	\textbf{最近距离}则要先看直线的位置：若 \(|C|>R\)，取二者较小值；
	若 \(|C|\le R\)，原直线与椭圆相交或相切，最近距离就是 \(0\)。

	\medskip

	\textbf{\textcolor{CExplanation}{几何本质（法向投影半宽）}}

	设直线 \(l:Ax+By+C=0\) 的法向量为
	\[
		\vec n=(A,B).
	\]
	则
	\[
		R=\sqrt{a^2A^2+b^2B^2}
	\]
	表示线性表达式 \(Ax+By\) 在椭圆上的最大绝对值，即
	\[
		\max_{(x,y)\in E}|Ax+By|=R.
	\]

	若把法向量单位化，则椭圆在直线法向方向上的真实投影半宽为
	\[
		\frac{R}{\sqrt{A^2+B^2}}.
	\]

	因此，\(R\) 本身是未归一化法向量 \((A,B)\) 对应的支撑函数值，而
	\[
		\frac{R}{\sqrt{A^2+B^2}}
	\]
	才是距离意义下的法向投影半宽。

	\medskip

	\textbf{\textcolor{CExplanation}{代数意义（统一运算）}}

	利用柯西不等式：
	\[
		(Aa\cos\theta+Bb\sin\theta)^2
		\le
		(a^2A^2+b^2B^2)(\cos^2\theta+\sin^2\theta)
		=
		R^2.
	\]

	这说明
	\[
		|Aa\cos\theta+Bb\sin\theta|\le R.
	\]

	并且当
	\[
		(\cos\theta,\sin\theta)
	\]
	与
	\[
		(aA,bB)
	\]
	同向或反向时，等号可以取到。因此
	\[
		M=Aa\cos\theta+Bb\sin\theta
	\]
	的值域确为
	\[
		[-R,R].
	\]

	\medskip

	\textbf{\textcolor{CExplanation}{考点价值（应试效率）}}

	\begin{itemize}[leftmargin=*, itemsep=0.5em, topsep=0.4em]
		\item \textbf{考法一（直接计算最值）：}
		      给出具体椭圆和直线，直接代入公式得最值，省去分类讨论的时间。

		\item \textbf{考法二（逆向求参）：}
		      已知距离最值的范围，利用不等式关系反求椭圆参数 \(a,b\) 或直线参数 \(A,B,C\)。
	\end{itemize}

	\medskip

	\textbf{\textcolor{CExplanation}{顿悟点}}

	\textbf{\(R\) 是解决此类问题的枢纽。只要算出 \(R\)，并比较 \(|C|\) 与 \(R\) 的大小，最值和位置关系就能同时确定。}

	\medskip

	\textbf{\textcolor{CExplanation}{使用场景}}

	\begin{itemize}[leftmargin=*, itemsep=0.5em, topsep=0.4em]
		\item \textbf{场景一（求距离最值）：}
		      已知椭圆和直线方程，直接套公式。

		\item \textbf{场景二（位置关系判断）：}
		      比较 \(|C|\) 与 \(R\) 的大小，判断相离、相切或相交。

		\item \textbf{场景三（反求参数）：}
		      已知距离最值，利用
		      \[
			      d_{\max}=\frac{|C|+R}{\sqrt{A^2+B^2}}
		      \]
		      反求椭圆参数或直线参数。
	\end{itemize}

\end{explanationbox}